To see whether more folding still lowers the S-event negative fraction, we ran
$8\times4\times1$ with the cubic map and the sum over sets, which are the defaults, after PR~\#305.
These runs were done on the cluster, with the same seed and number of events as the
corresponding $2\times2\times1$ and $4\times4\times1$ runs. Table~\ref{tab:floor} compares
$\rS$ for the three foldings. The CPU of these runs is not comparable with the desktop runs
and is not quoted.

For each run we also give, in parentheses, the expected contribution of the virt events,
$N_{\mathrm{virt}}/(2N_{\mathrm S})$. It assumes that half of them are negative, as for a residual
$V-A$ symmetric around zero; $N_{\mathrm{virt}}$ is taken from the generation log, since the
events themselves do not record whether they are virt events.

\begin{table}[h]
\centering\small
\resizebox{\textwidth}{!}{%
\begin{tabular}{lccccccc}
\toprule
 & \multicolumn{3}{c}{$\rS$ (virt-event part)} & & \multicolumn{3}{c}{$r$}\\
process & $2\times2\times1$ & $4\times4\times1$ & $8\times4\times1$ & $\rH$ & 2,2,1 & 4,4,1 & 8,4,1\\
\midrule
\csname @@input\endcsname floor_rows.txt
\bottomrule
\end{tabular}}
\caption{S-event negative fraction with the cubic map and the sum over FKS sets, for three
foldings. In parentheses: expected contribution of the virt events. $\rH$ and the total
$r$ are for the $8\times4\times1$ runs and for each folding respectively.}
\label{tab:floor}
\end{table}

The floor is reached in every process, at the latest at $4\times4\times1$:
\begin{itemize}
\item \textbf{$pp\to t\bar t$.} $\rS$ equals the virt-event contribution at every folding
($0.010$, $0.015$, $0.011$). All the negative S events are virt events, and the floor is
reached already at $2\times2\times1$. Folding cannot do more here; a better virtual
surrogate (the dipole fits of the virt-trick study) could.
\item \textbf{$pp\to e^+e^-j$.} $\rS$ goes from $0.122$ ($2\times2\times1$) to $0.086$
($4\times4\times1$) and $0.094\pm0.004$ ($8\times4\times1$). It stops at about $0.09$, of which
$0.02$ comes from virt events. The remaining $0.07$ is consistent with the floor of
about $0.03$--$0.07$ found at fixed Born points with 256 points (table~\ref{tab:dumpr}).
\item \textbf{$e^+e^-\to b\bar b b\bar b$.} $0.183$, $0.142$, then $0.150\pm0.005$: the floor is
at about $0.145$, of which about $0.06$ comes from the virt events, which here are 8\% of all
events.
\item \textbf{$pp\to Zjj$.} Running (result pending).

\end{itemize}
The part of the floor that is not due to virt events corresponds to Born points where the
radiation integral of the S terms is negative, or so steep near the dead zones and cuts
that no affordable number of points resolves its sign. In all four processes the total $r$
no longer moves between $4\times4\times1$ and $8\times4\times1$: beyond
$2\times2\times1$ or $4\times4\times1$, extra folding only costs CPU.
